This paper studies how local labor market tightness and welfare benefit levels at the time of receiving protection shape refugees’ labor market integration and location choices. Exploiting a sequence of reforms that generated quasi-random variation in benefit levels and labor market tightness across initial placements, we find that tighter local labor markets raise employment, while more generous benefits slightly reduce it. These effects persist for about 2.5 years, but both shocks have lasting impacts on internal migration. Moreover, the negative employment effect of benefits is larger in tighter labor markets. Neither shock changes the probability of leaving Austria altogether.
